\chapter[How Empires Made a World]{How Classical Empires Made All under Heaven and a World}
\section{A Bronze Coin}
Begin with a small object: a bronze coin.

Round, square-holed, polished at the edge, it bears two Chinese characters, \textit{ban liang}, half a \textit{liang}. Such coins circulated in China for centuries, from Qin over two millennia ago into early Han. Almost weightless in your palm, it nevertheless places an entire imperial system there.

Before it, the seven Warring States had separate currencies. Qi's coins resembled small knives, Yan's half a spade, Chu's wrinkled bronze faces called ant-nose money. A Zhao trader carrying knife money into Chu resembled someone carrying arcade tokens to an airport: nobody accepted it. Every cross-border purchase began with argument over what money was worth and how to exchange it.

Qin conquered the six states and ordered one form everywhere: round, square-holed, weighing half a liang. A Shu trader could buy Yan horses with the same currency that paid a Xianyang clerk and collected taxes throughout the realm. The hole mattered too: stringing coins facilitated counting, transport, and storage.

The coin answers three questions: who issues it, the court; what it is worth, a unified standard; who may use it, everyone under heaven. The hand determining these is our protagonist, the classical empire.

China was not alone. Around similar times, Persians used consistently fine gold darics, Romans silver bearing emperors' heads, Mauryans punch-marked square silver. Why did civilisations separated by seas and deserts, scarcely meeting, make similar things? What shared problem were they solving? Behind this coin lies technology for manufacturing all under heaven and a world.

\section{A Morning in Qianling County}
Come into a scene.

Dawn shortly after Qin's unification in 221 BCE. Qianling County, Dongting Commandery, in mountainous northwestern present-day Hunan. The Youshui curves beside town; mist has not dispersed.

Mortars sound first as convicts begin pounding rice. The county office opens. A clerk enters carrying trimmed wooden slips, sits, grinds ink. Today, like yesterday, he registers.

The entries are daunting. A household's locality, head's name, age, height---Qin registered height because it determined adulthood and service liability---members, adult males, ranks, last year's land tax. Beside him another checks postal records: date and hour of a document from Youyang, courier's name, lateness punishable under law. In a corner, convict-work slips record rice pounded and wall-building days. Above these come commandery summaries: annual rent, corvee levies, grain stocks.

This is not invention. Two thousand years later, archaeologists recovered discarded county records from an old well: over thirty-eight thousand written wooden slips, the Liye Qin documents. They even include a multiplication table, a clerk's calculating reference.

Stand here and notice the easily missed fact: nobody wears armour or holds a blade, yet this quiet room is the empire itself. Qin's final conquests took ten years, but conquest only began the work. Clerks behind slips truly made the realm, turning households into lines, fields into numbers, mountain roads into timetables. Far away in Xianyang, the emperor needed know no Qianling resident to count recruitable men and removable grain.

\section[All under Heaven Had to Be Made]{All under Heaven Was Made---and There Lies the Problem}
Set out the conflict before concluding.

The effectiveness was extraordinary. Unified writing let Chu officials read Qi documents; standard axle widths made imperial-road travel swift and steady; common measures aligned lengths and weights. Roads, Great Wall, Lingqu Canal moved grain from Lingnan to northern fronts. Two millennia later, Chinese characters, common measures, and life within the concept China follow these tracks.

The costs were extraordinary too. Who built walls, roads, palaces, and the Mount Li tomb? People selected line by line from offices like Qianling. Adult men owed yearly labour and frontier service. The Qin-ending rebellion at Dazexiang began with nine hundred conscripted farmers stranded by rain en route to the frontier.

The real question emerges, not historical trivia but a continuing dispute:

\textbf{Does a great order encompassing everyone protect or dispossess?}

Supporters say without it interstate wars continue, currencies clash, dams go unbuilt, bandits unchecked. Opponents say ordinary lives and sweat underpin it: some enjoy peace under heaven, others pay only with life.

Do not dismiss this as Qin brutality. The same recipe---roads, taxation, armies, law, currency, registration---appeared around similar periods in almost all great empires: Persia, Rome, Mauryan India, even a variant among the unwriting Inca. A tyrant's whim would not recur so widely. It resembles a shared human predicament: organising millions of strangers into one order. Every solution carries a bill.

Before continuing, ask whether in Qianling you would prefer registering others or being registered.

\section{Emperors, Conscripts, and the Conquered}
Hear multiple voices, giving each its full case.

\textbf{First: order's builders.} Strengthen emperors' and ministers' reasoning. Unification sounds naturally good to us and thus often escapes justification, which does thought no favour.

Four reasons at least. End war: over two Warring States centuries brought near-annual major battles; at Changping, Qin killed hundreds of thousands of surrendered Zhao soldiers. Better one authority than seven mutual killers. Convenience: common writing, measures, and money sharply lower transaction costs, opening trade, commands, and technical diffusion. Major works: irrigation, frontier walls, imperial roads exceed small states' capacity and materially benefit nearby farmers. Predictability: published written law and rule-bound officials improve on local nobles' whims.

The case is strong. Early Han recovery brought population growth and lower grain prices. During Rome's peace, Mediterranean piracy was cleared and shipping flourished. Large-scale order supplied real dividends.

\textbf{Second: those on the bill.} The other coin-face holds the payers.

In the \textit{Records of the Grand Historian}'s ``Hereditary House of Chen She,'' rain strands Chen Sheng, Wu Guang, and nine hundred conscripts at Dazexiang in Qi County, threatening their deadline. Sima Qian says they privately reasoned that flight and rebellion both meant death, so why not fight? Chen Sheng then shouted words still burning two millennia later:

\begin{quote}
``Are kings, nobles, generals, and ministers born of a special stock?''
\end{quote}
Notice the target: not order itself but supposedly predetermined places within it. Registered commoner households meant birth and ledger entries almost settled a life. Order gave the realm unity while pinning most people to a cell.

\textbf{Third: the conquered.} Often absent from imperial narratives. First-century Roman historian Tacitus described his father-in-law's British campaigns through a fierce speech attributed to Calgacus, a Scottish tribal leader approaching defeat: plunder, slaughter, theft they call rule; they make a desert and call it peace.

A Roman wrote this, imagining an enemy's speech, evidence of perceptive ears within Rome. Yet its point is real: Roman peace centred on the Mediterranean while frontiers bled. For Britons and Jews, tax collectors, crosses, and slave markets defined the imperial world.

\textbf{Fourth: the remorseful victor.} Rarer still, someone winning then stopping. Third-century BCE Mauryan ruler Ashoka conquered coastal Kalinga, then carved remorse in stones across his empire. The thirteenth edict describes devastation and announces conquest through dhamma instead of force: compassionate, tolerant, dutiful principles winning hearts. A ruler publicly admitting a victorious war wrong is exceptionally rare. Empires never contain one voice; even universal order's inventors can wake troubled at night.

\textbf{Fifth: those on the bottom step.} Look lower, at Roman slavery. Captives, piracy victims, and children of enslaved parents worked mines, estates, homes; literate people served as scribes and tutors. Roman law allowed manumission, even citizenship for freed people, whose children became full citizens. Not proof of kindness but of a status ladder: enslaved at bottom, freed and foreigners above, citizens with voting, legal, and corporal-punishment privileges atop. Qin had twenty ranks, promotion for battlefield merit, possibilities of redemption through service for convicts and enslaved people. Imperial classification was always intricate.

Pause at a misleading counterexample. In 212 CE, Caracalla granted citizenship to almost all free imperial inhabitants. Apparently equality's triumph. Yet when everyone became a citizen, citizenship lost value, and new divisions followed: higher and lower status, lighter and heavier punishment for identical offences. Remove an old ladder and a new one soon appears. Judge universal reforms by the cells redrawn in practice, not merely decree wording.

\section[Thinking Bridge: The Bill for Order]{Thinking Bridge: Draw Up the Bill for Any Great Order}
Can a portable tool prevent competing voices from carrying us away? Four plain questions. Whenever you hear a grand proposal---imperial road, city metro, school rule, platform feature---fill out its bill:

\textbf{First: What does it provide?} Honestly acknowledge benefits. Roads are fast, common currency convenient, Roman seas safer. Criticism denying benefits is as lazy as praise overlooking costs.

\textbf{Second: Who receives them?} Who uses roads? Roman armies and dispatches first, traders incidentally; Qin imperial carriages and official documents, ordinary travellers beside them. Who sails safe seas? Rich merchants and slave traders. One project pays different dividends by position.

\textbf{Third: Who pays?} Who carries wall bricks, supplies army grain, sends sons to frontiers? Rome leased tax collection to merchant groups, whose excess collections became profit after fixed payments, burdening provinces further. Whose pockets supplied the excess? Remember Qianling: each benefit has an expenditure entry somewhere, bearing actual names.

\textbf{Fourth: Who sets the rules?} Whose measures prevail, whose convenience shapes law? Chen Sheng's challenge asks why the same people always decide.

These questions often yield no simple good-or-bad verdict but a distribution: benefits here, costs there, rule-making hands between. More useful than declaring any dynasty's merits greater or smaller than faults.

Try school uniforms. Provision: easier management, less comparison. Beneficiaries: school, parents? Payers: whose individuality, whose money? Makers: do pupils vote? The tool crosses eras.

\section[Rain and Stone: Two Primary Texts]{A Rainstorm and a Stone: Reading Two Primary Texts}
Read two short passages, one from losers, one from a victor.

First, the already encountered ``Hereditary House of Chen She,'' where Sima Qian records calculations after the missed deadline:

\begin{quote}
``Flight now means death; launching a great undertaking also means death. Since death is equal, may we die for the state?''
\end{quote}
Broadly, if flight and uprising both kill, why not risk life in public affairs? Then comes the challenge to the special birth of kings and ministers.

An evidential qualification is necessary. Excavated Shuihudi Qin legal slips, a clerk's burial copies, punish ordinary corvee lateness through escalating fines and allow rain exemption, not universal beheading. How then explain the account's death penalty? Scholars differ: stricter military law for garrison conscripts, emergency late-Qin severity, or Chen Sheng's mobilising rhetoric. Written statute, temporary orders, rallying slogans, and enforcement are never identical. Before trusting ironclad evidence, identify its kind.

Second, the victor describes failure. Ashoka's thirteenth rock edict appeared beside roads throughout the empire for travellers. Broadly:

\begin{quote}
In his eighth regnal year, the Beloved of the Gods conquered Kalinga. One hundred and fifty thousand were deported, one hundred thousand killed, many times more died. \,\ldots\, The Beloved of the Gods feels deep remorse, for conquest of an unsubdued land entails killing, death, and deportation, weighing heavily upon him.
\end{quote}
Weigh both. People in a ledger tear it up; its author declares some entries crimes. About a century and thousands of kilometres apart, unaware of each other, they ask whether order's costs can be seen and acknowledged. Chen Sheng uses blades; Ashoka inscribes remorse. More often there is neither rebellion nor monument, only unread costs quietly resting in expenditure columns.

\section{All Roads Open---for Whom?}
Now the great fact: classical empires seemingly agreed to build roads and posts, mint coins, establish law, and register households.

Persia's Royal Road ran over two thousand kilometres from Sardis to Susa, with relay stations along it. Herodotus describes couriers unstoppable by snow, rain, heat, or night. Roman roads radiated to provinces, leaving today's proverb that all roads lead to Rome; milestones gave distances to the next city. Qin had imperial and straight roads, Han posts reached the Western Regions. Even the Andean Inca, without writing, money, or wheels, built tens of thousands of kilometres of mountain routes, transmitting numbers through quipu knots and maintaining bridges through mita labour levies.

Every companion component appears. Tax: Han poll levies and land rent, Persian fixed provincial tribute, Roman tax farmers. Law: Roman commoners forced publication on twelve bronze tables in the square; Qin specified daily convict rice-pounding quantities. Armies: Rome's provincial auxiliaries earned citizenship through completed service, turning armies into Roman-making machines; Qin recruited directly through registers. Other routes existed: India's later Gupta dynasty, roughly fourth--sixth centuries, granted land and taxation to dependants and temples rather than reaching each field. Local autonomy remained under an imperial umbrella.

Two movements ran against roads' apparent freedom. Empires also uprooted entire populations into required cells. Qin moved 120,000 powerful households to Xianyang, strengthening the capital and weakening local roots. After Meng Tian pushed back the Xiongnu, new counties received interior farmers for frontier cultivation and defence. Han Western Regions soldiers held hoe and spear. Roman veterans settled provincial colonies carrying law, language, loyalty: living imperial infrastructure. Ashoka's 150,000 deportees remind us that forced movement was war's spoil; displaced people often survive in history as totals alone.

Frontiers were therefore broad zones, not neat lines. The Great Wall and Roman Britain's walls managed movement through taxes, checks, admission, exclusion. Trade and marriage continued: Chinese silk and steppe horses exchanged at passes; guards and outside herders often knew one another best. Correct a common mistake: a wall that fails to block every mounted nomad was not necessarily wasted. Its tasks were delay, forcing slower crossings; warning, beacons outrunning cavalry; management, channelling trade and people through taxable, inspectable gates. Under these criteria, the Great Wall often worked normally. Judge infrastructure by intended function.

Why such convergence? Interpretations now compete. Three follow:

\textbf{First: control.} Roads, posts, registers, common money chiefly move orders, soldiers, taxes. Evidence: Roman military road specifications, Persian couriers carrying official papers, Qin imperial-road privileges. This states motives plainly without romance. Its limit: unexplained spillovers, as merchants, monks, and migrants used facilities and genuine economic integration occurred.

\textbf{Second: integration.} Standards lower transaction costs; large markets grow and yield imperial taxes. Evidence: Mediterranean trade under Roman peace, widening Silk Roads after Zhang Qian, currency turning distant bargaining into calculation. This explains imperial revenue. Its limit: trade prosperity does not mean universal benefit. Goods included humans; slave trading thrived on peaceful seas.

\textbf{Third: display.} Roads, monuments, standard vessels, magnificent capitals perform power, making distant residents see the centre and habituating obedience. Ashoka's roadside edicts, Persian trilingual cliff records of suppressed revolts, Qin measure-standardisation decrees cast on vessels: limited practical use, enormous visibility. Politics is stage as well as accounting. Its limit: if all is performance, where do actual taxes and troops come from? Theatre cannot replace granaries.

Each has truth; their overlap matters. One infrastructure is whip, circulatory system, and poster simultaneously. Use the bill: not whether a road is good or bad, but who benefits and pays under each function.

\section[Further Challenge: Making Society Legible]{Further Challenge: The State Must First See Society Clearly}
Now a theory linking all components. In \textit{Seeing Like a State}, twentieth-century political scientist James C. Scott makes a plain observation: \textbf{traditional states confront society as an unreadable fog}.

Local people know nicknames, who borrowed whose rice, who actually cultivates a field. Outsiders cannot see through it. Taxes, recruitment, judgement fail in the fog. The state's first task is therefore neither road-building nor legislation but \textbf{making society legible}: measuring and numbering fields into maps, fixing names and ages into registers, standardising measures into accounts, unifying different scripts into common communication.

Qianling's office now appears a machine converting fog into tables. Height, rank, rent, delivery times translate messy local life into numbers readable in Xianyang. Scott adds an overlooked example: many communities originally lacked fixed surnames, changing names by context. Registration gradually imposed stable surnames. A name lasting from birth certificate to gravestone owes partly to tradition, partly to archives.

The theory illuminates Qin's script, axle, and measure standards; Roman registration categories of citizen, Latin, foreigner; Persian tribute totals. It also clarifies solutions to universal order versus local variation. Qin demands one table. Persia is looser: Cyrus's cylinder claims return of peoples deported to Babylon and restored temples; Ezra records permission for Jews to rebuild Jerusalem's temple. Pay taxes, do not rebel, remain yourself. Ashoka communicates one dhamma in local writing, including Kharoshthi, Greek, Aramaic in western regions: universal ideals in local dress. Gupta incorporates local diversity through dependants and temples. One realm, four blueprints.

The sharper half follows: \textbf{legibility cuts both ways and often damages things}. Simplified-away knowledge---which slope needs fallow, old grievances shaping judgement, crop rotation, social credit arrangements---is survival wisdom accumulated over centuries. Pruning society into tables may prune wisdom too. Qin's urgency becomes overforceful legibility: forcing a realm into tables within decades until the tables align and society breaks.

The reverse appears in Scott's \textit{The Art of Not Being Governed}. If visibility entails taxes, service, battle, people evade it through mountains, marshes, forests, beyond ledgers. His Southeast Asian highland argument suggests supposed backwardness may reflect ancestors deliberately leaving states, not missing civilisation. Chinese displaced people and fugitives, Roman barbarians, unregistered Persian tribes become similar frontier answers: \textbf{we will not take your roads or enter your registers}.

Weigh the theory and its limits. Not every state act reduces to legibility: armies, faith, personal ambition matter. Clarity is not inherently bad; without registration, no school rolls or relief lists. Theory is a searchlight illuminating much while darkening what lies outside its beam.

\section{The Empire in Your Pocket}
This ancient technology operates upon you now.

Your lifelong identity number is a legible code, finer than a Qin entry specifying locality and height but continuous in function. Education offices count school-age children through enrolment systems. Map apps turn cities into live tables of congestion. A Qianling clerk seeing navigation would consider wooden slips a preliminary draft.

Daily dividends arrive: a nationally uniform entrance examination, usable-everywhere renminbi, spontaneous high-speed rail trips, readable road signs. Costs arrive too, more gently: systems register your information, platforms and institutions read you beyond any emperor's dreams. Big data extends legibility to steps and location.

Universal order versus local difference also persists. Mandarin enables nationwide communication while dialects quietly recede, nonstandard yet carrying grandmother's storytelling voice. Common curricula measure educational equality but compress local rhythms. Even delivery-address fields---province, city, district, street, number---would impress Persian couriers. Time too is unified. All China uses Beijing time despite western sunrise over two hours later, making Urumqi's ten o'clock workday seem late-rising. One time zone writes order across the sky. SIM-card and account real-name checks share Qianling's simple rationale: find the person when something happens.

Frontiers and migration echo through visas, passports, border checks, negotiations between ledgers and those beyond them. Every move from county town to provincial capital and farther repeats the choice: enter larger order or remain outside its reach.

\section[Your Turn: How Visible Will You Be?]{Your Turn: Would You Let Yourself Be Seen Completely?}
Return to Qianling's office at dawn.

A clerk records name, age, height. The state now knows you: taxes recorded, service counted, offences investigated. Protection also begins: land has a deed, grievances an address, famine relief a named entry.

We now hold four billing questions, three road explanations, legibility's searchlight, Ashoka's stone, Chen Sheng's rain, Tacitus's desert called peace.

The question is yours:

\textbf{If an order's ability to see you is exactly proportional to its ability to protect you and serve itself---clearer means safer, but fewer hiding places---how far would you let it see?}

Answer with reasons. Reweigh the other faces: registration also means relief lists; square-holed coins open trade; unified writing accompanies dialect loss; Dazexiang's rain stands opposite winters the Great Wall truly defended. Remember those refusing roads and registers: theirs too is an answer.

Two thousand years ago the registering clerk himself appeared in another book. Today your name appears in many. Perhaps the difference begins now: you know what those books are, how they arose, and that they have never been free.
